\documentclass{resume}

\usepackage[left=0.5in,top=0.4in,right=0.5in,bottom=0.4in]{geometry}

\name{Raghav Sareen}
\address{(650) 660-7750 $\diamond$ Pleasanton, CA}
\address{\href{mailto:raghavsareen2006@gmail.com}{raghavsareen2006@gmail.com} $\diamond$
\href{https://www.linkedin.com/in/raghavsareen2006/}{linkedin.com/in/raghavsareen2006} $\diamond$
\href{https://github.com/RaghavSareen2006}{github.com/RaghavSareen2006}}

\begin{document}

\begin{rSection}{Objective}
Computer Engineering student pursuing Firmware and AI Engineering opportunities, with experience spanning embedded C/C++, STM32 motor control, edge-AI computer vision, LLM evaluation, agent testing, Linux systems, and hardware-software integration.
\end{rSection}

\begin{rSection}{Education}
\textbf{Bachelor of Science in Computer Engineering}, San Jose State University \hfill Expected May 2028\\
GPA: 3.776/4.000\\
\textit{Selected Coursework:} Programming Concepts, Object-Oriented Programming, Algorithms and Data Structures, Assembly Language Programming, Electronics, Circuit Analysis, Introductory Electrical Engineering Lab, Engineering Statistics, Differential Equations and Linear Algebra
\end{rSection}

\begin{rSection}{Skills}
\begin{tabular}{@{} >{\bfseries}l @{\hspace{2ex}} p{0.73\textwidth}}
Programming & C, C++ \\
Firmware \& Embedded & STM32F103, libHAL, GPIO, CAN, PWM, I\textsuperscript{2}C, PID control, BLDC motors, limit switches \\
AI \& Computer Vision & Ultralytics YOLO, TensorRT, OpenCV, prompt engineering, LLM evaluation, agentic workflows \\
Systems \& DevOps & NVIDIA Jetson Orin Nano, ARM64 Linux, Docker Compose, Tailscale, Git/GitHub \\
AI \& Product Tools & OpenClaw, Claude Code, ChatGPT, Jira, Kanban, Scrum, PRDs, epics, features, user stories \\
\end{tabular}
\end{rSection}

\begin{rSection}{Technical Experience}

\textbf{Firmware Team Member, ALICANTO Rover} \hfill Sep 2024 -- May 2026\\
SJSU Robotics \hfill \textit{San Jose, CA}
\begin{itemize}
    \itemsep -3pt {}
    \item Programmed and tested C++ firmware for the ALICANTO rover arm subsystem in preparation for the 2027 University Rover Challenge; committed changes that were reviewed and merged into the shared team codebase.
    \item Validated three STM32F103 MicroMod and motor assemblies controlling the wrist, shoulder, and elbow, verifying motor response and component behavior with libHAL.
    \item Updated and tuned PID control functions to reduce arm jitter and improve stability during movement.
    \item Collaborated with the firmware lead and sublead to implement arm homing for BLDC motors using one limit switch and hands-on GPIO, CAN, and PWM communication.
\end{itemize}

\textbf{AI Agent Testing Volunteer} \hfill Jun 2026\\
Social Experiments \hfill \textit{Remote}
\begin{itemize}
    \itemsep -3pt {}
    \item Designed 60 functional test cases and executed 37 functional tests plus 13 end-to-end workflow scenarios for an AI Slack agent.
    \item Tested Jira integration, contract drafting, document ingestion, redlining, approvals, role permissions, duplicate handling, error recovery, and response-time behavior.
    \item Built a structured test workbook containing priorities, preconditions, prompts, test steps, expected and actual results, automation candidates, and defect notes.
    \item Identified failures or errors in 28 scenarios, including duplicate responses, long timeouts, incorrect task routing, missing attachments, and incomplete workflow transitions.
    \item Reported results to the product team, contributing to improvements across most identified issues; used Jira and Agile practices and learned to structure PRDs with epics, features, and user stories.
\end{itemize}

\end{rSection}

\newpage
\begin{rSection}{Projects}

\textbf{AI-Assisted Home Security Arm Turret} \hfill Jul 2026 -- Present\\
\textit{NVIDIA Jetson Orin Nano, xArm, YOLO, TensorRT, OpenCV, Docker Compose, Tailscale}
\begin{itemize}
    \itemsep -3pt {}
    \item Conceived and assembled a customizable home-security system by integrating a USB camera and a commercial xArm kit with a Jetson Orin Nano; configured two servo motors for pan and tilt control.
    \item Used OpenClaw and Claude Code for AI-assisted implementation, then personally deployed, configured, troubleshot, and validated the complete hardware-software system for continuous 24/7 operation at home.
    \item Tested person and face detection, target verification and reacquisition, automatic room scanning, dog tracking, manual hold, digital zoom, photo capture, Telegram alerts, GPIO status lights, and safe servo limits.
    \item Deployed GPU inference and web services with Docker Compose and built private remote access through Tailscale; validated an authenticated dashboard with livestream controls and real-time telemetry.
\end{itemize}

\textbf{Jetson Home Lab and Personal AI Infrastructure} \hfill Jul 2026 -- Present\\
\textit{Jetson Orin Nano, ARM64 Linux, OpenClaw, Docker Compose, Tailscale, Home Assistant}
\begin{itemize}
    \item Configured the Jetson as a multi-service home server running OpenClaw, Home Assistant, Music Assistant, Jellyfin, Portainer, and computer-vision containers.
    \item Managed persistent volumes, host networking, container health, logs, service restarts, and private HTTPS access through Tailscale.
    \item Diagnosed IoT connection failures, camera conflicts, container crash loops, high system load, swap pressure, and USB HID availability; configured exit-node forwarding and persistent NAT/firewall behavior.
    \item Created a sanitized Home Assistant configuration repository with Docker Compose files, restore documentation, secret exclusions, and credential-pattern checks.
\end{itemize}

\textbf{The Efficacy of ChatGPT in Mathematics} \hfill May 2023 -- Sep 2023\\
\textit{Polygence Research Program}
\begin{itemize}
    \itemsep -3pt {}
    \item Designed and conducted an experiment evaluating ChatGPT across 107 algebra problems from a standardized textbook dataset.
    \item Measured 81\% initial answer accuracy, compared two follow-up strategies, and analyzed failure patterns involving graphs, rational expressions, and mathematical formatting.
    \item Evaluated LaTeX formatting and math-teacher role prompting as mitigation strategies; authored a research paper and presentation with guidance from a research mentor.
    \item Won Top Asynchronous Pitch among more than 300 presentations at the Polygence Symposium of Rising Scholars.
\end{itemize}

\end{rSection}

\begin{rSection}{Leadership \& Activities}

\textbf{Social Media Coordinator, Indian Students Organization} \hfill Aug 2025 -- Present\\
Created two promotional reels, participated in four to five event videos, attended weekly planning meetings, and supported promotion and event operations for a program that sold 75 tickets.

\textbf{Program Member, MESA Engineering Program} \hfill Sep 2024 -- Present\\
Volunteered 12+ hours as an event floater for Science Extravaganza, supporting hands-on STEM workshops serving more than 350 middle-school students.

\end{rSection}

\begin{rSection}{Awards \& Certifications}
\textbf{Awards:} 2026 President's Scholar, San Jose State University; Top Asynchronous Pitch, Polygence Symposium of Rising Scholars; AP Scholar; California State Seal of Biliteracy\\
\textbf{Certifications:} Command Line in Linux and Getting Started with Linux Terminal (Coursera, 2026); Learn Git and GitHub in One Day (Udemy, 2026); Claude Code in Action (Anthropic, 2026)
\end{rSection}

\begin{rSection}{Languages}
English (Full Professional Proficiency) $\diamond$ Hindi (Native or Bilingual Proficiency) $\diamond$ Spanish (Limited Working Proficiency)
\end{rSection}

\end{document}
